\abstractEN % Do NOT modify this line

Existing digital learning platforms often lack effective adaptive personalisation and content oriented towards civic literacy and urban sustainability --- gaps that contribute to high dropout rates in informal education contexts.

This work proposes \textit{Play4Change}, an adaptive and gamified learning platform designed for citizens, developed within the U!REKA mission~\cite{ureka2024} and aligned with the United Nations Sustainable Development Goals 4, 11 and~13~\cite{un2015sdg}. The central contribution is an adaptive path reuse mechanism: whenever an artificial intelligence agent generates an alternative learning path for a citizen experiencing difficulties, that path is indexed by vector similarity and automatically reused for citizens with similar difficulty profiles, reducing inference cost and progressively accumulating collective pedagogical knowledge.

The platform integrates a cross-platform mobile client (Android and iOS), a web administration portal, an AI agent based on \textit{LangChain4j} with the \textit{Mistral} model and semantic retrieval via \textit{pgvector}, and a \textit{stateless} server implemented in \textit{Spring Boot}. Learning is structured into atomic micro-competencies, recognised through badges. The system was implemented as a functional prototype and validated against the requirements defined in collaboration with the U!REKA mission.
